% !TeX program = lualatex
% ================================================================
% Polish Parallel Text Version of FractionsQuickQuestions
%
% Generated by the server using the following settings:
%
%   language            Polish (pol)
%   text direction      left to right
%   font                Noto Serif
%   fallback font       Noto Serif
%   built from          latex/quickQuestions/fractionsQuickQuestions.tex
%   vocabulary key      latex/quickQuestions/fractionsQuickQuestions/L2/pol/fractionsQuickQuestions_polishKey.tex
%   tier 3 vocabulary   green   RGB 0 128 0
%   English text        black   RGB 0 0 0
%   Polish text         purple  RGB 112 48 160
%   layout macros       version 3
%
% Please don't edit any information in this comment block - the server
% compares it against the current settings and rebuilds this file when
% they differ, so changes here would be overwritten. However, feel free
% to make updates to the LaTeX below if you want to edit it.
% ================================================================

\documentclass[aspectratio=169]{beamer}
% ================================================================
% Copied from the English sheet this was translated from, so that
% anything it sets up for itself still works here
% ================================================================
\usepackage{amsmath}
\usepackage{amssymb}

% ================================================================
% Quick Questions deck.
%
% One question per slide, no answers: the teacher jumps to a topic
% from the contents slide and picks questions for mini whiteboards.
% Within each topic the ten questions get gradually harder.
% ================================================================

\setbeamertemplate{navigation symbols}{}
\setbeamertemplate{footline}{}

% Topic divider.
\newcommand{\qtopic}[1]{%
  \section{#1}%
  \begin{frame}[plain]
    \vfill
    \begin{center}
      {\usebeamercolor[fg]{structure}\Huge\bfseries #1}
    \end{center}
    \vfill
  \end{frame}%
}

% A question slide: the question on its own, centred, nothing else on
% the slide. The optional argument drops the type size for the wordier
% questions, so nothing runs off the slide.
\newcommand{\qq}[2][\Huge]{%
  \begin{frame}
    \vfill
    \begin{center}
      \begin{minipage}{0.92\textwidth}
        \centering #1 #2
      \end{minipage}
    \end{center}
    \vfill
  \end{frame}%
}

% Contents-slide bullets, colour-coded by difficulty band:
% green for the first four topics, orange for the next four,
% red for the last three.
\setbeamertemplate{section in toc}{%
  \ifnum\inserttocsectionnumber<5
    \textcolor{green!55!black}{$\bullet$}%
  \else\ifnum\inserttocsectionnumber<9
    \textcolor{orange}{$\bullet$}%
  \else
    \textcolor{red}{$\bullet$}%
  \fi\fi
  ~\inserttocsection%
}

\title{Fractions: Quick Questions}
\subtitle{Mini whiteboard questions}
\author{}
\date{}

% Characters this language's font has not got are borrowed from another,
% rather than being silently dropped - see docs/translated-sheet-latex.md
\directlua{%
  if luaotfload and luaotfload.add_fallback then
    luaotfload.add_fallback("eallatinfallback",
      { "Noto Serif:mode=harf;" })
  else
    tex.error("this luaotfload is too old to register a fallback font")
  end
}
\usepackage[english]{babel}
\babelprovide[import]{polish}
\babelfont[polish]{rm}[Renderer=Harfbuzz, BoldFont={Noto Serif}, ItalicFont={Noto Serif}, BoldItalicFont={Noto Serif}, RawFeature={fallback=eallatinfallback}]{Noto Serif}
\babelfont[polish]{sf}[Renderer=Harfbuzz, BoldFont={Noto Serif}, ItalicFont={Noto Serif}, BoldItalicFont={Noto Serif}, RawFeature={fallback=eallatinfallback}]{Noto Serif}
\newcommand{\ealtext}[1]{\foreignlanguage{polish}{#1}}
\newcommand{\ealtextblock}[1]{{\raggedright\foreignlanguage{polish}{#1}\par}}
% ================================================================
% EAL layout helpers
% ================================================================
\usepackage{xcolor}

\definecolor{ealkeycolour}{RGB}{0,128,0}
\definecolor{ealenglish}{RGB}{0,0,0}
\definecolor{ealtwocolour}{RGB}{112,48,160}

\newsavebox{\ealboxen}
\newsavebox{\ealboxtr}
\newlength{\ealglosswd}
\newlength{\ealglossraise}
\setlength{\ealglossraise}{1.15em}

% a tier 3 word, in English and in translation alike
\newcommand{\ealkey}[1]{\textcolor{ealkeycolour}{#1}}
\newcommand{\ealkeytr}[1]{\textcolor{ealkeycolour}{\ealtext{#1}}}

% One translated word above one English word. Built from kernel
% primitives, never a tabular, which would break wherever the sheet
% loads array, booktabs or tabularx. The box is as wide as the wider of
% the two, so a long translation pushes its neighbours aside instead of
% overlapping them, and it claims its full height so a table row or a
% TikZ node makes room for it.
\newcommand{\ealgloss}[2]{%
  \sbox{\ealboxen}{\ealkey{#1}}%
  \sbox{\ealboxtr}{\small\textcolor{ealtwocolour}{\ealtext{#2}}}%
  \setlength{\ealglosswd}{\wd\ealboxen}%
  \ifdim\wd\ealboxtr>\ealglosswd\setlength{\ealglosswd}{\wd\ealboxtr}\fi
  \leavevmode
  \makebox[\ealglosswd][c]{%
    \makebox[0pt][c]{\usebox{\ealboxen}}%
    \makebox[0pt][c]{\raisebox{\ealglossraise}{\usebox{\ealboxtr}}}%
  }%
}

% opens up line spacing around a block of glosses, so the raised
% translations sit clear of the line above
\newenvironment{ealglossed}{\par\linespread{2.1}\selectfont\ignorespaces}{\par}

% a whole translated sentence above its English counterpart. block level,
% so it does not belong inside a tabular cell
\newcommand{\ealpara}[2]{%
  \par\addvspace{0.5\baselineskip}%
  {\color{ealtwocolour}\ealtextblock{#1}}%
  \penalty200
  {\color{ealenglish}#2\par}%
  \addvspace{0.5\baselineskip}%
}

\begin{document}

\frame{\titlepage}

\begin{frame}{Contents}
  \small
  \tableofcontents
\end{frame}

%================================================================
\qtopic{Simplifying fractions}
%================================================================

\qq{\ealpara{\ealkeytr{Uprość} $\dfrac{4}{8}$}{\ealkey{Simplify} $\dfrac{4}{8}$}}

\qq{\ealpara{\ealkeytr{Uprość} $\dfrac{6}{9}$}{\ealkey{Simplify} $\dfrac{6}{9}$}}

\qq{\ealpara{\ealkeytr{Uprość} $\dfrac{12}{20}$}{\ealkey{Simplify} $\dfrac{12}{20}$}}

\qq{\ealpara{\ealkeytr{Uprość} $\dfrac{24}{36}$}{\ealkey{Simplify} $\dfrac{24}{36}$}}

\qq{\ealpara{\ealkeytr{Uprość} $\dfrac{45}{60}$}{\ealkey{Simplify} $\dfrac{45}{60}$}}

\qq{\ealpara{\ealkeytr{Uprość} $\dfrac{42}{70}$}{\ealkey{Simplify} $\dfrac{42}{70}$}}

\qq{\ealpara{\ealkeytr{Uprość} $\dfrac{84}{144}$}{\ealkey{Simplify} $\dfrac{84}{144}$}}

\qq{\ealpara{\ealkeytr{Uprość} $\dfrac{91}{117}$}{\ealkey{Simplify} $\dfrac{91}{117}$}}

\qq{\ealpara{\ealkeytr{Uprość} $\dfrac{315}{420}$}{\ealkey{Simplify} $\dfrac{315}{420}$}}

\qq[\LARGE]{\ealpara{%
\ealkeytr{Uprość} $\dfrac{1001}{1309}$\\[0.4em]
\normalsize (Wskazówka: obie liczby są \ealkeytr{wielokrotnościami} liczb $7$ i $11$.)%
}{%
\ealkey{Simplify} $\dfrac{1001}{1309}$\\[0.4em]
\normalsize (Hint: both numbers are \ealkey{multiples} of $7$ and of $11$.)%
}}

%================================================================
\qtopic{Ordering fractions}
%================================================================

\qq[\LARGE]{\ealpara{%
Ustaw te liczby od najmniejszej do największej:\\[0.6em]
$\dfrac{1}{4}\qquad \dfrac{1}{2}\qquad \dfrac{1}{8}$%
}{%
Put these in order, smallest first:\\[0.6em]
$\dfrac{1}{4}\qquad \dfrac{1}{2}\qquad \dfrac{1}{8}$%
}}

\qq[\LARGE]{\ealpara{%
Ustaw te liczby od najmniejszej do największej:\\[0.6em]
$\dfrac{3}{7}\qquad \dfrac{6}{7}\qquad \dfrac{2}{7}$%
}{%
Put these in order, smallest first:\\[0.6em]
$\dfrac{3}{7}\qquad \dfrac{6}{7}\qquad \dfrac{2}{7}$%
}}

\qq[\LARGE]{\ealpara{%
Ustaw te liczby od najmniejszej do największej:\\[0.6em]
$\dfrac{1}{2}\qquad \dfrac{3}{4}\qquad \dfrac{5}{8}$%
}{%
Put these in order, smallest first:\\[0.6em]
$\dfrac{1}{2}\qquad \dfrac{3}{4}\qquad \dfrac{5}{8}$%
}}

\qq[\LARGE]{\ealpara{%
Ustaw te liczby od najmniejszej do największej:\\[0.6em]
$\dfrac{2}{3}\qquad \dfrac{5}{6}\qquad \dfrac{3}{4}$%
}{%
Put these in order, smallest first:\\[0.6em]
$\dfrac{2}{3}\qquad \dfrac{5}{6}\qquad \dfrac{3}{4}$%
}}

\qq[\LARGE]{\ealpara{%
Ustaw te liczby od najmniejszej do największej:\\[0.6em]
$\dfrac{3}{5}\qquad \dfrac{7}{10}\qquad \dfrac{13}{20}$%
}{%
Put these in order, smallest first:\\[0.6em]
$\dfrac{3}{5}\qquad \dfrac{7}{10}\qquad \dfrac{13}{20}$%
}}

\qq[\LARGE]{\ealpara{%
Ustaw te liczby od najmniejszej do największej:\\[0.6em]
$\dfrac{5}{9}\qquad \dfrac{2}{3}\qquad \dfrac{7}{12}$%
}{%
Put these in order, smallest first:\\[0.6em]
$\dfrac{5}{9}\qquad \dfrac{2}{3}\qquad \dfrac{7}{12}$%
}}

\qq[\LARGE]{\ealpara{%
Ustaw te liczby od największej do najmniejszej:\\[0.6em]
$\dfrac{4}{5}\qquad \dfrac{7}{9}\qquad \dfrac{11}{15}$%
}{%
Put these in order, \textbf{largest} first:\\[0.6em]
$\dfrac{4}{5}\qquad \dfrac{7}{9}\qquad \dfrac{11}{15}$%
}}

\qq[\LARGE]{\ealpara{%
Ustaw te liczby od najmniejszej do największej:\\[0.6em]
$\dfrac{5}{8}\qquad \dfrac{7}{11}\qquad \dfrac{9}{14}$%
}{%
Put these in order, smallest first:\\[0.6em]
$\dfrac{5}{8}\qquad \dfrac{7}{11}\qquad \dfrac{9}{14}$%
}}

\qq[\LARGE]{\ealpara{%
Która liczba jest bliżej $1$: $\dfrac{8}{9}$ czy $\dfrac{11}{12}$?\\[0.4em]
\normalsize Skąd wiesz?%
}{%
Which is closer to $1$: $\dfrac{8}{9}$ or $\dfrac{11}{12}$?\\[0.4em]
\normalsize How do you know?%
}}

\qq[\LARGE]{\ealpara{%
Napisz \ealkeytr{ułamek}, który leży między\\[0.6em]
$\dfrac{5}{8}$ \quad a \quad $\dfrac{2}{3}$%
}{%
Write a \ealkey{fraction} that lies between\\[0.6em]
$\dfrac{5}{8}$ \quad and \quad $\dfrac{2}{3}$%
}}

%================================================================
\qtopic{Mixed numbers and improper fractions}
%================================================================

\qq{\ealpara{%
Zapisz $1\dfrac{1}{2}$ w postaci \ealkeytr{ułamka niewłaściwego}%
}{%
Write $1\dfrac{1}{2}$ as an \ealkey{improper fraction}%
}}

\qq{\ealpara{%
Zapisz $2\dfrac{3}{4}$ w postaci \ealkeytr{ułamka niewłaściwego}%
}{%
Write $2\dfrac{3}{4}$ as an \ealkey{improper fraction}%
}}

\qq{\ealpara{%
Zapisz $\dfrac{7}{2}$ w postaci \ealkeytr{liczby mieszanej}%
}{%
Write $\dfrac{7}{2}$ as a \ealkey{mixed number}%
}}

\qq{\ealpara{%
Zapisz $4\dfrac{2}{5}$ w postaci \ealkeytr{ułamka niewłaściwego}%
}{%
Write $4\dfrac{2}{5}$ as an \ealkey{improper fraction}%
}}

\qq{\ealpara{%
Zapisz $\dfrac{23}{4}$ w postaci \ealkeytr{liczby mieszanej}%
}{%
Write $\dfrac{23}{4}$ as a \ealkey{mixed number}%
}}

\qq{\ealpara{%
Zapisz $5\dfrac{7}{8}$ w postaci \ealkeytr{ułamka niewłaściwego}%
}{%
Write $5\dfrac{7}{8}$ as an \ealkey{improper fraction}%
}}

\qq{\ealpara{%
Zapisz $\dfrac{100}{7}$ w postaci \ealkeytr{liczby mieszanej}%
}{%
Write $\dfrac{100}{7}$ as a \ealkey{mixed number}%
}}

\qq{\ealpara{%
Zapisz $12\dfrac{5}{9}$ w postaci \ealkeytr{ułamka niewłaściwego}%
}{%
Write $12\dfrac{5}{9}$ as an \ealkey{improper fraction}%
}}

\qq{\ealpara{%
Zapisz $-3\dfrac{2}{5}$ w postaci \ealkeytr{ułamka niewłaściwego}%
}{%
Write $-3\dfrac{2}{5}$ as an \ealkey{improper fraction}%
}}

\qq[\LARGE]{\ealpara{%
\ealkeytr{Liczba mieszana} $6\dfrac{a}{7}$ jest taka sama jak $\dfrac{45}{7}$.\\[0.6em]
Znajdź $a$.%
}{%
The \ealkey{mixed number} $6\dfrac{a}{7}$ is the same as $\dfrac{45}{7}$.\\[0.6em]
Find $a$.%
}}

%================================================================
\qtopic{Adding and subtracting fractions}
%================================================================

\qq{$\dfrac{1}{2}+\dfrac{1}{4}$}

\qq{$\dfrac{2}{3}+\dfrac{1}{6}$}

\qq{$\dfrac{3}{4}-\dfrac{1}{8}$}

\qq{$\dfrac{1}{3}+\dfrac{1}{4}$}

\qq{$\dfrac{3}{5}-\dfrac{1}{4}$}

\qq{$\dfrac{5}{6}+\dfrac{3}{8}$}

\qq{$\dfrac{7}{10}-\dfrac{2}{15}$}

\qq{$\dfrac{2}{3}+\dfrac{3}{4}+\dfrac{1}{6}$}

\qq{$\dfrac{5}{12}-\dfrac{7}{18}$}

\qq[\LARGE]{\ealpara{%
Co należy dodać do $\dfrac{5}{9}$,\\[0.4em]
aby otrzymać $\dfrac{7}{8}$?%
}{%
What must be added to $\dfrac{5}{9}$\\[0.4em]
to make $\dfrac{7}{8}$?%
}}

%================================================================
\qtopic{Converting fractions, decimals and percentages}
%================================================================

\qq[\LARGE]{\ealpara{%
Zapisz $\dfrac{1}{2}$ jako \ealkeytr{ułamek dziesiętny} i jako \ealkeytr{procent}%
}{%
Write $\dfrac{1}{2}$ as a \ealkey{decimal} and as a \ealkey{percentage}%
}}

\qq[\LARGE]{\ealpara{%
Zapisz $0.25$ jako \ealkeytr{ułamek} i jako \ealkeytr{procent}%
}{%
Write $0.25$ as a \ealkey{fraction} and as a \ealkey{percentage}%
}}

\qq[\LARGE]{\ealpara{%
Zapisz $30\%$ jako \ealkeytr{ułamek dziesiętny} i jako \ealkeytr{ułamek} w \ealkeytr{najprostszej postaci}%
}{%
Write $30\%$ as a \ealkey{decimal} and as a \ealkey{fraction} in its \ealkey{simplest form}%
}}

\qq[\LARGE]{\ealpara{%
Zapisz $\dfrac{3}{4}$ jako \ealkeytr{procent}, a $0.6$ jako \ealkeytr{procent}%
}{%
Write $\dfrac{3}{4}$ as a \ealkey{percentage}, and $0.6$ as a \ealkey{percentage}%
}}

\qq[\LARGE]{\ealpara{%
Zapisz $\dfrac{7}{20}$ jako \ealkeytr{ułamek dziesiętny} i jako \ealkeytr{procent}%
}{%
Write $\dfrac{7}{20}$ as a \ealkey{decimal} and as a \ealkey{percentage}%
}}

\qq[\LARGE]{\ealpara{%
Zapisz $0.08$ jako \ealkeytr{procent} i jako \ealkeytr{ułamek} w \ealkeytr{najprostszej postaci}%
}{%
Write $0.08$ as a \ealkey{percentage} and as a \ealkey{fraction} in its \ealkey{simplest form}%
}}

\qq[\LARGE]{\ealpara{%
Zapisz $\dfrac{3}{8}$ jako \ealkeytr{ułamek dziesiętny} i jako \ealkeytr{procent}%
}{%
Write $\dfrac{3}{8}$ as a \ealkey{decimal} and as a \ealkey{percentage}%
}}

\qq[\LARGE]{\ealpara{%
Zapisz $12.5\%$ jako \ealkeytr{ułamek} w \ealkeytr{najprostszej postaci}%
}{%
Write $12.5\%$ as a \ealkey{fraction} in its \ealkey{simplest form}%
}}

\qq[\LARGE]{\ealpara{%
Zapisz $\dfrac{5}{6}$ jako \ealkeytr{ułamek dziesiętny} oraz jako \ealkeytr{procent} z dokładnością do $1$ \ealkeytr{miejsca dziesiętnego}%
}{%
Write $\dfrac{5}{6}$ as a \ealkey{decimal}, and as a \ealkey{percentage} to $1$ \ealkey{decimal place}%
}}

\qq[\LARGE]{\ealpara{%
Zapisz $137.5\%$ jako \ealkeytr{liczbę mieszaną} w \ealkeytr{najprostszej postaci}%
}{%
Write $137.5\%$ as a \ealkey{mixed number} in its \ealkey{simplest form}%
}}

%================================================================
\qtopic{Ordering fractions, decimals and percentages}
%================================================================

\qq[\LARGE]{\ealpara{%
Ustaw te liczby od najmniejszej do największej:\\[0.6em]
$0.5\qquad \dfrac{1}{4}\qquad 60\%$%
}{%
Put these in order, smallest first:\\[0.6em]
$0.5\qquad \dfrac{1}{4}\qquad 60\%$%
}}

\qq[\LARGE]{\ealpara{%
Ustaw te liczby od najmniejszej do największej:\\[0.6em]
$0.7\qquad \dfrac{3}{4}\qquad 70\%$%
}{%
Put these in order, smallest first:\\[0.6em]
$0.7\qquad \dfrac{3}{4}\qquad 70\%$%
}}

\qq[\LARGE]{\ealpara{%
Ustaw te liczby od najmniejszej do największej:\\[0.6em]
$\dfrac{2}{5}\qquad 45\%\qquad 0.35$%
}{%
Put these in order, smallest first:\\[0.6em]
$\dfrac{2}{5}\qquad 45\%\qquad 0.35$%
}}

\qq[\LARGE]{\ealpara{%
Ustaw te liczby od najmniejszej do największej:\\[0.6em]
$0.62\qquad \dfrac{5}{8}\qquad 63\%$%
}{%
Put these in order, smallest first:\\[0.6em]
$0.62\qquad \dfrac{5}{8}\qquad 63\%$%
}}

\qq[\LARGE]{\ealpara{%
Ustaw te liczby od największej do najmniejszej:\\[0.6em]
$\dfrac{7}{8}\qquad 0.87\qquad 88\%$%
}{%
Put these in order, \textbf{largest} first:\\[0.6em]
$\dfrac{7}{8}\qquad 0.87\qquad 88\%$%
}}

\qq[\LARGE]{\ealpara{%
Ustaw te liczby od najmniejszej do największej:\\[0.6em]
$\dfrac{1}{3}\qquad 0.3\qquad 33\%$%
}{%
Put these in order, smallest first:\\[0.6em]
$\dfrac{1}{3}\qquad 0.3\qquad 33\%$%
}}

\qq[\LARGE]{\ealpara{%
Ustaw te liczby od najmniejszej do największej:\\[0.6em]
$0.045\qquad 4.5\%\qquad \dfrac{1}{20}$%
}{%
Put these in order, smallest first:\\[0.6em]
$0.045\qquad 4.5\%\qquad \dfrac{1}{20}$%
}}

\qq[\LARGE]{\ealpara{%
Ustaw te liczby od najmniejszej do największej:\\[0.6em]
$\dfrac{5}{6}\qquad 83\%\qquad 0.835$%
}{%
Put these in order, smallest first:\\[0.6em]
$\dfrac{5}{6}\qquad 83\%\qquad 0.835$%
}}

\qq[\LARGE]{\ealpara{%
Ustaw te liczby od najmniejszej do największej:\\[0.6em]
$\dfrac{2}{7}\qquad 0.28\qquad 29\%$%
}{%
Put these in order, smallest first:\\[0.6em]
$\dfrac{2}{7}\qquad 0.28\qquad 29\%$%
}}

\qq[\LARGE]{\ealpara{%
Ustaw te liczby od najmniejszej do największej:\\[0.6em]
$\dfrac{3}{8}\qquad 0.38\qquad 37.5\%\qquad \dfrac{3}{7}$\\[0.5em]
\normalsize Co zauważasz?%
}{%
Put these in order, smallest first:\\[0.6em]
$\dfrac{3}{8}\qquad 0.38\qquad 37.5\%\qquad \dfrac{3}{7}$\\[0.5em]
\normalsize What do you notice?%
}}

%================================================================
\qtopic{Fractions of amounts (no calculator)}
%================================================================

\qq{\ealpara{$\dfrac{1}{2}$ z $40$}{$\dfrac{1}{2}$ of $40$}}

\qq{\ealpara{$\dfrac{1}{4}$ z $60$}{$\dfrac{1}{4}$ of $60$}}

\qq{\ealpara{$\dfrac{1}{5}$ z $45$}{$\dfrac{1}{5}$ of $45$}}

\qq{\ealpara{$\dfrac{3}{4}$ z $40$}{$\dfrac{3}{4}$ of $40$}}

\qq{\ealpara{$\dfrac{2}{3}$ z $36$}{$\dfrac{2}{3}$ of $36$}}

\qq{\ealpara{$\dfrac{3}{5}$ z $45$}{$\dfrac{3}{5}$ of $45$}}

\qq{\ealpara{$\dfrac{5}{8}$ z $72$}{$\dfrac{5}{8}$ of $72$}}

\qq{\ealpara{$\dfrac{2}{7}$ z $350$}{$\dfrac{2}{7}$ of $350$}}

\qq[\LARGE]{\ealpara{%
$\dfrac{5}{6}$ pewnej liczby to $45$.\\[0.4em]
Jaka to liczba?%
}{%
$\dfrac{5}{6}$ of a number is $45$.\\[0.4em]
What is the number?%
}}

\qq[\LARGE]{\ealpara{%
$\dfrac{3}{8}$ z $\dfrac{2}{5}$ z $200$%
}{%
$\dfrac{3}{8}$ of $\dfrac{2}{5}$ of $200$%
}}

%================================================================
\qtopic{Fractions of amounts (calculator)}
%================================================================

\qq{\ealpara{$\dfrac{1}{4}$ z $386$}{$\dfrac{1}{4}$ of $386$}}

\qq{\ealpara{$\dfrac{3}{5}$ z $245$}{$\dfrac{3}{5}$ of $245$}}

\qq{\ealpara{$\dfrac{2}{7}$ z $483$}{$\dfrac{2}{7}$ of $483$}}

\qq{\ealpara{$\dfrac{3}{16}$ z $1024$}{$\dfrac{3}{16}$ of $1024$}}

\qq{\ealpara{$\dfrac{5}{9}$ z $738$}{$\dfrac{5}{9}$ of $738$}}

\qq[\LARGE]{\ealpara{%
$\dfrac{7}{8}$ z $\pounds 312.40$%
}{%
$\dfrac{7}{8}$ of $\pounds 312.40$%
}}

\qq{\ealpara{$\dfrac{11}{13}$ z $2405$}{$\dfrac{11}{13}$ of $2405$}}

\qq[\LARGE]{\ealpara{%
$\dfrac{4}{7}$ z $3.5$ kg.\\[0.4em]
Podaj odpowiedź w gramach.%
}{%
$\dfrac{4}{7}$ of $3.5$ kg.\\[0.4em]
Give your answer in grams.%
}}

\qq[\LARGE]{\ealpara{%
$\dfrac{5}{8}$ pewnej liczby to $215$.\\[0.4em]
Jaka to liczba?%
}{%
$\dfrac{5}{8}$ of a number is $215$.\\[0.4em]
What is the number?%
}}

\qq[\LARGE]{\ealpara{%
$\dfrac{7}{12}$ z czasu $4$ godzin i $48$ minut.\\[0.4em]
Podaj odpowiedź w godzinach i minutach.%
}{%
$\dfrac{7}{12}$ of $4$ hours $48$ minutes.\\[0.4em]
Give your answer in hours and minutes.%
}}

%================================================================
\qtopic{Four operations with mixed numbers and negatives}
%================================================================

\qq{$1\dfrac{1}{2}+2\dfrac{1}{4}$}

\qq{$3\dfrac{4}{5}-1\dfrac{2}{5}$}

\qq{$2\dfrac{1}{2}+\left(-1\dfrac{1}{4}\right)$}

\qq{$4\dfrac{1}{6}-2\dfrac{3}{4}$}

\qq{$2\dfrac{1}{2}\times 4$}

\qq{$-3\dfrac{1}{3}+1\dfrac{5}{6}$}

\qq{$1\dfrac{1}{2}\times 2\dfrac{2}{3}$}

\qq{$3\dfrac{3}{4}\div 1\dfrac{1}{2}$}

\qq{$-5\dfrac{1}{3}\div \dfrac{2}{3}$}

\qq[\LARGE]{$-2\dfrac{1}{4}\times\left(-1\dfrac{1}{3}\right)+1\dfrac{1}{2}$}

%================================================================
\qtopic{Word problems with mixed numbers}
%================================================================

\qq[\LARGE]{\ealpara{%
Przepis wymaga $1\dfrac{1}{2}$ szklanki mąki.\\[0.4em]
Alex robi podwójną porcję.\\[0.4em]
Ile mąki potrzeba?%
}{%
A recipe needs $1\dfrac{1}{2}$ cups of flour.\\[0.4em]
Alex makes a double batch.\\[0.4em]
How much flour is needed?%
}}

\qq[\LARGE]{\ealpara{%
Deska ma $3\dfrac{1}{4}$\,m długości.\\[0.4em]
Odcina się kawałek o długości $1\dfrac{1}{2}$\,m.\\[0.4em]
Ile deski zostanie?%
}{%
A plank is $3\dfrac{1}{4}$\,m long.\\[0.4em]
A piece $1\dfrac{1}{2}$\,m long is cut off.\\[0.4em]
How much plank is left?%
}}

\qq[\LARGE]{\ealpara{%
Priya ma dwie wstążki o długościach $2\dfrac{1}{3}$\,m i $1\dfrac{5}{6}$\,m.\\[0.4em]
Jaka jest ich łączna długość?%
}{%
Priya has two ribbons, $2\dfrac{1}{3}$\,m and $1\dfrac{5}{6}$\,m long.\\[0.4em]
What is their total length?%
}}

\qq[\LARGE]{\ealpara{%
Sam w poniedziałek przebiegł $2\dfrac{3}{4}$ mili, a we wtorek $3\dfrac{1}{2}$ mili.\\[0.4em]
O ile dalej przebiegł we wtorek?%
}{%
Sam runs $2\dfrac{3}{4}$ miles on Monday and $3\dfrac{1}{2}$ miles on Tuesday.\\[0.4em]
How much further did he run on Tuesday?%
}}

\qq[\LARGE]{\ealpara{%
Puszka zawiera $4\dfrac{1}{2}$ litra farby.\\[0.4em]
Wylewa się pięć dzbanków po $\dfrac{3}{4}$ litra.\\[0.4em]
Ile farby zostanie?%
}{%
A tin holds $4\dfrac{1}{2}$ litres of paint.\\[0.4em]
Five jugs of $\dfrac{3}{4}$ of a litre are poured out.\\[0.4em]
How much paint is left?%
}}

\qq[\LARGE]{\ealpara{%
Ściana ma $7\dfrac{1}{2}$\,m szerokości.\\[0.4em]
Panele mają $1\dfrac{1}{4}$\,m szerokości.\\[0.4em]
Ile paneli zmieści się wzdłuż ściany?%
}{%
A wall is $7\dfrac{1}{2}$\,m wide.\\[0.4em]
Panels are $1\dfrac{1}{4}$\,m wide.\\[0.4em]
How many panels fit across the wall?%
}}

\qq[\LARGE]{\ealpara{%
Butelka zawiera $1\dfrac{1}{3}$ litra wody.\\[0.4em]
Jo wypija $\dfrac{3}{8}$ tej wody.\\[0.4em]
Ile wody wypija Jo?%
}{%
A bottle holds $1\dfrac{1}{3}$ litres of water.\\[0.4em]
Jo drinks $\dfrac{3}{8}$ of it.\\[0.4em]
How much water does Jo drink?%
}}

\qq[\LARGE]{\ealpara{%
Podróż trwa $2\dfrac{1}{4}$ godziny.\\[0.4em]
Maria przejechała $\dfrac{2}{3}$ drogi.\\[0.4em]
Ile minut zostało?%
}{%
A journey takes $2\dfrac{1}{4}$ hours.\\[0.4em]
Maria has driven $\dfrac{2}{3}$ of the way.\\[0.4em]
How long is left, in minutes?%
}}

\qq[\LARGE]{\ealpara{%
Linę o długości $12\dfrac{1}{2}$\,m pocięto na kawałki o długości $1\dfrac{3}{4}$\,m.\\[0.4em]
Ile jest całych kawałków i ile liny zostanie?%
}{%
A rope $12\dfrac{1}{2}$\,m long is cut into pieces $1\dfrac{3}{4}$\,m long.\\[0.4em]
How many whole pieces are there, and how much rope is left over?%
}}

\qq[\large]{\ealpara{%
Przepis dla $4$ osób używa $2\dfrac{1}{2}$ szklanek ryżu i $1\dfrac{3}{4}$ litra bulionu.\\[0.5em]
Ile ryżu i bulionu potrzeba dla $10$ osób?%
}{%
A recipe for $4$ people uses $2\dfrac{1}{2}$ cups of rice and $1\dfrac{3}{4}$ litres of stock.\\[0.5em]
How much rice and stock are needed for $10$ people?%
}}

%================================================================
\qtopic{Recurring decimals to fractions}
%================================================================

\qq[\LARGE]{\ealpara{%
Użyj \ealkeytr{metody algebraicznej}, aby zapisać\\[0.5em]
$0.\dot{3}$ \quad jako \ealkeytr{ułamek}%
}{%
Use an \ealkey{algebraic method} to write\\[0.5em]
$0.\dot{3}$ \quad as a \ealkey{fraction}%
}}

\qq[\LARGE]{\ealpara{%
Użyj \ealkeytr{metody algebraicznej}, aby zapisać\\[0.5em]
$0.\dot{7}$ \quad jako \ealkeytr{ułamek}%
}{%
Use an \ealkey{algebraic method} to write\\[0.5em]
$0.\dot{7}$ \quad as a \ealkey{fraction}%
}}

\qq[\LARGE]{\ealpara{%
Użyj \ealkeytr{metody algebraicznej}, aby zapisać\\[0.5em]
$0.\dot{4}\dot{5}$ \quad jako \ealkeytr{ułamek}%
}{%
Use an \ealkey{algebraic method} to write\\[0.5em]
$0.\dot{4}\dot{5}$ \quad as a \ealkey{fraction}%
}}

\qq[\LARGE]{\ealpara{%
Użyj \ealkeytr{metody algebraicznej}, aby zapisać\\[0.5em]
$0.\dot{1}\dot{2}$ \quad jako \ealkeytr{ułamek} w \ealkeytr{najprostszej postaci}%
}{%
Use an \ealkey{algebraic method} to write\\[0.5em]
$0.\dot{1}\dot{2}$ \quad as a \ealkey{fraction} in its \ealkey{simplest form}%
}}

\qq[\LARGE]{\ealpara{%
Użyj \ealkeytr{metody algebraicznej}, aby zapisać\\[0.5em]
$0.\dot{2}\dot{7}$ \quad jako \ealkeytr{ułamek} w \ealkeytr{najprostszej postaci}%
}{%
Use an \ealkey{algebraic method} to write\\[0.5em]
$0.\dot{2}\dot{7}$ \quad as a \ealkey{fraction} in its \ealkey{simplest form}%
}}

\qq[\LARGE]{\ealpara{%
Użyj \ealkeytr{metody algebraicznej}, aby zapisać\\[0.5em]
$0.\dot{1}2\dot{3}$ \quad jako \ealkeytr{ułamek} w \ealkeytr{najprostszej postaci}%
}{%
Use an \ealkey{algebraic method} to write\\[0.5em]
$0.\dot{1}2\dot{3}$ \quad as a \ealkey{fraction} in its \ealkey{simplest form}%
}}

\qq[\LARGE]{\ealpara{%
Użyj \ealkeytr{metody algebraicznej}, aby zapisać\\[0.5em]
$0.1\dot{6}$ \quad jako \ealkeytr{ułamek} w \ealkeytr{najprostszej postaci}%
}{%
Use an \ealkey{algebraic method} to write\\[0.5em]
$0.1\dot{6}$ \quad as a \ealkey{fraction} in its \ealkey{simplest form}%
}}

\qq[\LARGE]{\ealpara{%
Użyj \ealkeytr{metody algebraicznej}, aby zapisać\\[0.5em]
$0.8\dot{3}$ \quad jako \ealkeytr{ułamek} w \ealkeytr{najprostszej postaci}%
}{%
Use an \ealkey{algebraic method} to write\\[0.5em]
$0.8\dot{3}$ \quad as a \ealkey{fraction} in its \ealkey{simplest form}%
}}

\qq[\LARGE]{\ealpara{%
Użyj \ealkeytr{metody algebraicznej}, aby zapisać\\[0.5em]
$1.\dot{2}\dot{7}$ \quad jako \ealkeytr{liczbę mieszaną} w \ealkeytr{najprostszej postaci}%
}{%
Use an \ealkey{algebraic method} to write\\[0.5em]
$1.\dot{2}\dot{7}$ \quad as a \ealkey{mixed number} in its \ealkey{simplest form}%
}}

\qq[\LARGE]{\ealpara{%
Użyj \ealkeytr{metody algebraicznej}, aby pokazać, że\\[0.5em]
$0.\dot{9}=1$%
}{%
Use an \ealkey{algebraic method} to show that\\[0.5em]
$0.\dot{9}=1$%
}}

\end{document}